\chapter{Poisson Infrastructure}

\section{Stable Lévy Measure}

\begin{defn}[Stable Lévy density]\label{def:stable_density}
\lean{stableLevyDensity}
\leanok
The symmetric $\gamma$-stable Lévy measure density is
$\nu_\gamma(dz) = c_\gamma |z|^{-1-\gamma}\,dz$ on $\mathbb{R} \setminus \{0\}$,
where $c_\gamma > 0$ depends on the stability index $\gamma \in (1,2)$.
This measure governs the jump intensity of a symmetric $\gamma$-stable Lévy process:
small jumps are frequent (the density diverges at the origin) while large jumps are rare
(the density decays polynomially). The singularity at the origin is strong enough that
$\int |z|^2\,\nu_\gamma(dz) = \infty$, which is why the process has infinite variance
and the energy space is necessarily Banach rather than Hilbert.
\end{defn}

\begin{theorem}[Density symmetry]\label{thm:density_sym}
\lean{stableLevyDensity_symmetric}
\leanok
\uses{def:stable_density}
The stable Lévy density is symmetric: $\nu_\gamma(-z) = \nu_\gamma(z)$ for all
$z \neq 0$. This follows immediately from the absolute value in the definition
$c_\gamma |z|^{-1-\gamma}$. Symmetry of $\nu_\gamma$ implies that the driving
process $L$ satisfies $L \stackrel{d}{=} -L$ (distributional symmetry),
which is the key ingredient in the even-annihilation argument underlying
the representability obstruction.
\end{theorem}

\begin{theorem}[Density nonnegativity]\label{thm:density_nonneg}
\lean{stableLevyDensity_nonneg}
\leanok
\uses{def:stable_density}
The stable Lévy density is nonneg when $c_\gamma \geq 0$: since
$|z|^{-1-\gamma} > 0$ for $z \neq 0$ and $c_\gamma > 0$, the product
$c_\gamma |z|^{-1-\gamma}$ is strictly positive away from the origin.
This ensures that $\nu_\gamma$ is a valid (non-negative) measure on
$\mathbb{R} \setminus \{0\}$, as required for a Lévy measure.
Nonnegativity is used throughout the moment computations to guarantee
that integrals against $\nu_\gamma$ have definite sign.
\end{theorem}

\section{Moment Computations}

\begin{theorem}[Large jump rate positive]\label{thm:large_jump_rate}
\lean{stable_large_jump_rate_pos}
\leanok
The rate of large jumps is $\lambda = \nu_\gamma(\{|z| > 1\}) = 2c_\gamma/\gamma > 0$.
This is computed by integrating the Lévy density over $\{|z| > 1\}$:
$\int_{|z|>1} c_\gamma |z|^{-1-\gamma}\,dz = 2c_\gamma \int_1^\infty z^{-1-\gamma}\,dz
= 2c_\gamma/\gamma$. Positivity of $\lambda$ is essential for the obstruction argument:
it ensures that the Poisson random variable $G_0 = \#\{|\Delta L_s| > 1, s \leq T\}$
has parameter $\lambda T > 0$, hence positive variance, which ultimately shows
that the centered witness $G_0 - \mathbb{E}[G_0]$ is nonzero.
\end{theorem}

\begin{theorem}[Small jump moment finite]\label{thm:small_jump_finite}
\lean{stable_small_jump_moment_finite}
\leanok
For $\alpha > \gamma$, the small-jump moment $\int_{|z|\leq 1} |z|^\alpha\,\nu_\gamma(dz)
= 2c_\gamma/(\alpha - \gamma)$ is finite and positive. The integral converges because
the exponent $\alpha - 1 - \gamma > -1$ at the origin makes the integrand integrable.
This finiteness is needed in the truncated measure construction: the second moment
$\int_{|z|\leq M} z^2\,\nu_\gamma(dz) < \infty$ (the case $\alpha = 2 > \gamma$)
ensures that compound Poisson approximations have finite variance, enabling the
$L^2$ isometry for truncated compensated integrals.
\end{theorem}

\begin{theorem}[Small jump moment diverges]\label{thm:small_jump_diverges}
\lean{stable_small_jump_moment_diverges}
\leanok
For $\alpha < \gamma$, the small-jump moment integral diverges: $\alpha - \gamma < 0$,
so $\int_{|z|\leq 1} |z|^\alpha\,\nu_\gamma(dz) = \infty$. In particular, setting
$\alpha = \gamma$ gives $\int |z|^\gamma\,\nu_\gamma(dz) = \infty$, confirming that
$L_t \notin L^\gamma(\Omega)$. This divergence is the fundamental reason why the
integrand space $H_L = L^p$ with $p < \gamma < 2$ cannot be a Hilbert space:
the $L^2$ norm of the jump-size variable $z$ under $\nu_\gamma$ is infinite,
forcing the entire framework into the Banach setting.
\end{theorem}

\section{Poisson Identities}

\begin{defn}[Poisson recurrence]\label{lem:poisson_recurrence}
\lean{poissonPMFReal_succ_mul}
\leanok
The Poisson PMF satisfies the recurrence relation $(k+1) \cdot P_r(k+1) = r \cdot P_r(k)$
for all $k \geq 0$, where $P_r(k) = e^{-r} r^k / k!$ is the Poisson probability mass
function with parameter $r > 0$. This identity is verified by direct algebraic manipulation
of the factorial ratio $(k+1)! / k! = k+1$. The recurrence is the key computational tool
for proving both the mean and second factorial moment of the Poisson distribution: it
allows reindexing sums $\sum k \cdot P_r(k)$ into sums of the form $r \sum P_r(k)$,
which reduce to $r$ by PMF normalization.
\end{defn}

\begin{theorem}[Poisson mean]\label{thm:poisson_mean}
\lean{poisson_mean_identity}
\leanok
\uses{lem:poisson_recurrence}
The expected value of a Poisson random variable with parameter $r$ equals $r$:
$\sum_{n=0}^{\infty} n \cdot P_r(n) = r$.
The proof uses the recurrence $(k+1)P_r(k+1) = r \cdot P_r(k)$ to reindex the series,
reducing to the PMF normalization $\sum P_r(n) = 1$ via \texttt{HasSum} manipulation.
This identity is a proved theorem in Lean (not axiomatized), and is used in the variance
computation and to establish the centering of the large-jump count witness
$G_0 - \mathbb{E}[G_0]$ that drives the representability obstruction.
\end{theorem}

\begin{theorem}[Poisson second factorial moment]\label{thm:poisson_second_fact}
\lean{poisson_second_factorial_hassum}
\leanok
\uses{lem:poisson_recurrence}
The second factorial moment of a Poisson random variable with parameter $r$ equals $r^2$:
$\sum_{n=0}^{\infty} n(n-1) P_r(n) = r^2$.
The proof applies the recurrence relation twice: first to remove one factor of $n$
(reducing to $r$ times the first moment sum), then again to evaluate the resulting sum.
Combined with the mean identity, this yields the variance via
$\mathrm{Var} = \mathbb{E}[N(N-1)] + \mathbb{E}[N] - (\mathbb{E}[N])^2 = r^2 + r - r^2 = r$.
\end{theorem}

\begin{theorem}[Poisson variance]\label{thm:poisson_variance}
\lean{poisson_variance_identity}
\leanok
\uses{thm:poisson_mean, thm:poisson_second_fact}
The variance of a Poisson random variable with parameter $r$ equals $r$:
$\mathrm{Var}(\mathrm{Poisson}(r)) = r$. The proof combines the mean identity
$\mathbb{E}[N] = r$ and the second factorial moment $\mathbb{E}[N(N-1)] = r^2$
via the algebraic identity $\mathrm{Var}(N) = \mathbb{E}[N^2] - (\mathbb{E}[N])^2
= \mathbb{E}[N(N-1)] + \mathbb{E}[N] - (\mathbb{E}[N])^2 = r^2 + r - r^2 = r$.
This is the critical link in the obstruction chain: positive variance
$\mathrm{Var}(G_0) = \lambda T > 0$ implies that the centered witness
$G_0 - \mathbb{E}[G_0]$ is nonzero in $L^2$.
\end{theorem}

\section{Quadratic Defect}

\begin{theorem}[Quadratic defect exact]\label{thm:quad_exact}
\lean{quadratic_defect_exact}
\leanok
The quadratic defect identity states that
$(V + Kz)^2 - V^2 - 2V(Kz) = (Kz)^2$ for all real numbers $V$, $K$, $z$.
This is a purely algebraic identity obtained by expanding the left-hand side:
$(V + Kz)^2 = V^2 + 2V(Kz) + (Kz)^2$, so the first two terms cancel.
The identity isolates the second-order remainder in the Taylor expansion of
$f(x) = x^2$ around $x = V$ with increment $Kz$, and serves as the prototype
for the jump defect in the Itô formula for compound Poisson processes.
\end{theorem}

\begin{theorem}[Quadratic defect sharp]\label{thm:quad_sharp}
\lean{quadratic_defect_sharp}
\leanok
\uses{thm:quad_exact}
For the quadratic function $f(x) = x^2$, the jump remainder
$g_f(V,K,z) := f(V + Kz) - f(V) - f'(V) \cdot Kz$ equals exactly $(Kz)^2$.
This sharpens the quadratic defect identity by identifying the remainder with
the Itô correction term. The result shows that for quadratic functionals,
the Leibniz defect $\Gamma_L(F,F)$ is governed entirely by the sum of squared
jumps $\sum (\Delta L_s)^2$---the quadratic variation, which is a.s.\ infinite
for stable processes. This illustrates concretely why the classical Itô formula
(which uses quadratic variation) breaks down for $\gamma$-stable processes.
\end{theorem}

\section{Canonical Poisson Construction}

\begin{defn}[Concrete stable jump count]\label{def:concrete_jump}
\lean{ConcreteStableJumpCount}
\leanok
\uses{thm:large_jump_rate}
Given the stable parameters $(c_\gamma, \gamma, T)$, this constructs the canonical
Poisson probability space $(\mathbb{N}, \mathrm{poissonMeasure}(\lambda T))$ where
$\lambda = 2c_\gamma/\gamma$ is the large-jump rate. The identity function
$G_0(n) = n$ on this space serves as the concrete realization of the large-jump
count $\#\{s \in [0,T] : |\Delta L_s| > 1\}$. This construction is important
because it provides an explicit, machine-checkable probability space on which
the centered witness for the representability obstruction lives,
avoiding the need for an abstract existence argument.
\end{defn}

\begin{theorem}[$G_0$ has Poisson distribution]\label{thm:G0_poisson}
\lean{ConcreteStableJumpCount.G₀_poisson}
\leanok
\uses{def:concrete_jump}
The random variable $G_0$ on the canonical Poisson space has distribution
$\mathrm{Poisson}(\lambda T)$: $\mathrm{Measure.map}\;G_0\;\mu = \mathrm{poissonMeasure}(\lambda T)$.
The proof is trivial by $\mathrm{Measure.map\_id}$, since $G_0$ is the identity function
on $\mathbb{N}$ and the base measure is already $\mathrm{poissonMeasure}(\lambda T)$.
This seemingly tautological result is nonetheless essential in the formalization:
it provides the type-theoretic bridge between the abstract Poisson distribution
and the concrete random variable used in the obstruction argument.
\end{theorem}


\chapter{Lévy--Itô Framework}

\section{Truncated Measures}

\begin{defn}[Truncated stable measure]\label{def:trunc_measure}
\lean{TruncatedStableMeasure}
\leanok
The truncated stable measure is the restriction of $\nu_\gamma$ to $\{|z| \leq M\}$
for a fixed truncation level $M > 0$. This truncation removes the large-jump tail
of the Lévy measure, producing a finite measure
$\nu_\gamma^M := \nu_\gamma|_{\{|z|\leq M\}}$ with finite moments of all orders.
The truncated measure is the key technical device that enables compound Poisson
approximation: the process with Lévy measure $\nu_\gamma^M$ is a compound Poisson
process (finite activity), for which the Itô formula reduces to a finite telescoping sum.
The full stable integral is recovered by sending $M \to \infty$.
\end{defn}

\begin{theorem}[Truncated second moment positive]\label{thm:trunc_second}
\lean{TruncatedStableMeasure.second_moment_pos}
\leanok
\uses{def:trunc_measure}
The truncated second moment $\int_{|z|\leq M} z^2\,\nu_\gamma(dz) > 0$ is strictly
positive for any $M > 0$. This follows from the explicit computation
$\int_{|z|\leq M} z^2\,\nu_\gamma(dz) = 2c_\gamma M^{2-\gamma}/(2-\gamma)$
with $c_\gamma > 0$ and $2 - \gamma > 0$ (since $\gamma < 2$).
Positivity of this moment is needed for the $L^2$ isometry of the truncated
compensated integral $\delta_L^M$: the isometry constant involves
$\int z^2\,\nu_\gamma^M(dz)$, which must be nonzero to give a genuine norm.
\end{theorem}

\section{Finite Telescoping}

\begin{theorem}[Finite telescope]\label{thm:telescope}
\lean{finite_telescope}
\leanok
The finite telescoping identity states that
$f(y_n) - f(y_0) = \sum_{i=0}^{n-1} [f(y_{i+1}) - f(y_i)]$ for any function $f$
and finite sequence $y_0, \ldots, y_n$. The proof uses Mathlib's
\texttt{Finset.sum\_range\_sub}, which formalizes telescoping for finitely many terms.
This identity replaces stochastic integration in the compound Poisson setting:
since a compound Poisson process has finitely many jumps in $[0,T]$, the Itô
formula reduces to a finite sum of jump increments, requiring no measure-theoretic
integration theory. This makes the compound Poisson Itô formula fully formalizable
with elementary Lean machinery.
\end{theorem}

\begin{theorem}[Jump linear plus remainder]\label{thm:jump_decomp}
\lean{jump_linear_plus_remainder}
\leanok
Each jump increment decomposes as
$f(v+z) - f(v) = f'(v)z + [f(v+z) - f(v) - f'(v)z]$: a linear part $f'(v)z$
(which contributes to the stochastic integral $\int f'(L_{t-})\,dL_t$) plus a
nonlinear remainder (the jump correction). This is a trivial algebraic identity
but it isolates the two contributions that appear in the Lévy--Itô decomposition:
the first term is odd in $z$ (detectable by the one-parameter derivative $D_L$),
while the remainder involves the full nonlinear response of $f$ to each jump
and contributes to the Leibniz defect $\Gamma_L$.
\end{theorem}

\begin{theorem}[Product telescope]\label{thm:product_telescope}
\lean{product_telescope}
\leanok
\uses{thm:telescope}
The finite telescope applied to the product $h(x) = a(x) \cdot b(x)$ yields
the Itô product rule for compound Poisson processes: the increment
$(ab)(y_{i+1}) - (ab)(y_i)$ decomposes into a ``Leibniz part''
$a(y_i)[b(y_{i+1}) - b(y_i)] + b(y_i)[a(y_{i+1}) - a(y_i)]$ plus a
``defect part'' $[a(y_{i+1}) - a(y_i)][b(y_{i+1}) - b(y_i)]$.
Summing over all jumps via the telescope identity gives the product rule
with jump defect (Theorem~B) in the compound Poisson case.
This algebraic decomposition is the finite-dimensional prototype of the
infinite-dimensional Leibniz defect formula.
\end{theorem}

\section{Compound Poisson Integral}

\begin{defn}[Compound Poisson integral]\label{def:cpi}
\lean{CompoundPoissonIntegral}
\leanok
The compound Poisson integral is the finite random sum
$\sum_{i=1}^{N} h(s_i, z_i)$ minus its compensator
$\mathbb{E}[\sum_i h(s_i, z_i)]$, where $(s_i, z_i)$ are the jump times and sizes
of a compound Poisson process with Lévy measure $\nu_\gamma^M$ (restricted to
$\{|z| \leq M\}$). Since the truncated Lévy measure has finite total mass,
there are a.s.\ finitely many jumps, making the sum well-defined pointwise.
This object serves as the $\varepsilon$-level approximation $I_1^{\varepsilon,M}(h)$
in the construction of the compensated integral, and its linearity
(\texttt{random\_sum\_linear}) is inherited by the limit.
\end{defn}

\begin{theorem}[Random sum linear]\label{thm:random_sum_linear}
\lean{CompoundPoissonIntegral.random_sum_linear}
\leanok
\uses{def:cpi}
The compound Poisson integral is linear:
$\sum(c \cdot h_1 + h_2) = c \cdot \sum h_1 + \sum h_2$ for any constant $c$ and
integrands $h_1, h_2$. This follows from the linearity of finite sums and the
linearity of expectation for the compensator. Linearity of the finite-sum approximation
is the starting point for proving linearity of the full compensated integral $I_1^M(h)$:
since the limit of linear maps (in $L^2$) is linear, the constructed integral inherits
this property from its approximations.
\end{theorem}

\section{Compensated Integral Construction}

\begin{defn}[Constructed compensated integral]\label{def:cci}
\lean{ConstructedCompensatedIntegral}
\leanok
\uses{def:cpi}
The compensated integral $I_1^M(h) = \lim_{\varepsilon \to 0} I_1^{\varepsilon,M}(h)$
is constructed as the $L^2$ limit of compound Poisson finite sums as the lower truncation
$\varepsilon \to 0^+$. The $L^2$ Cauchy property of the approximations is guaranteed by
the estimate $\|I_1^{\varepsilon_1,M}(h) - I_1^{\varepsilon_2,M}(h)\|_2^2
\leq C \cdot \varepsilon^{2-\gamma}/(2-\gamma) \to 0$. This construction avoids
axiomatizing the stochastic integral: centering and isometry are \emph{derived}
from the corresponding properties of the finite-sum approximations via
\texttt{tendsto\_nhds\_unique}, making them proved theorems rather than assumptions.
\end{defn}

\begin{theorem}[$I_1$ centering from approximation]\label{thm:I1_centered}
\lean{ConstructedCompensatedIntegral.I₁_centered_from_approx}
\leanok
\uses{def:cci}
The constructed compensated integral is centered: $\mathbb{E}[I_1^M(h)] = 0$.
The proof proceeds by passing centering through the limit: each finite-sum
approximation $I_1^{\varepsilon,M}(h)$ is centered (by construction, as a compensated
sum), and since $I_1^{\varepsilon,M}(h) \to I_1^M(h)$ in $L^2$ (hence in $L^1$),
the expectation $\mathbb{E}[I_1^M(h)] = \lim \mathbb{E}[I_1^{\varepsilon,M}(h)] = 0$
by \texttt{tendsto\_nhds\_unique}. This derived centering feeds into the abstract
framework's hypothesis (H1): $\mathbb{E}[\delta_L(u)] = 0$ for all integrands $u$.
\end{theorem}

\begin{theorem}[$I_1$ isometry from approximation]\label{thm:I1_isometry}
\lean{ConstructedCompensatedIntegral.I₁_isometry_from_approx}
\leanok
\uses{def:cci}
The constructed compensated integral satisfies the $L^2$ isometry:
$\mathbb{E}[|I_1^M(h)|^2] = \|h\|_{L^2(dt \otimes \nu_\gamma^M)}^2$.
As with centering, this is derived from the isometry of the finite-sum approximations:
each $I_1^{\varepsilon,M}(h)$ satisfies an exact $L^2$ isometry (a standard property
of compensated Poisson integrals on finite measures), and the isometry identity is
preserved in the $L^2$ limit. This provides the norm equivalence needed for the
truncated integral and feeds into the BDG-type bound for $\delta_L^M$.
\end{theorem}

\begin{defn}[Bridge: construction to axioms]\label{def:toAxioms}
\lean{ConstructedCompensatedIntegral.toAxioms}
\leanok
\uses{thm:I1_centered, thm:I1_isometry}
This bridge packages the constructed compensated integral (with its derived centering
and isometry properties) into the axiomatic interface expected by the Lévy--Itô
decomposition layer. The bridge converts proved properties of the $L^2$-limit
construction into the structure fields required by \texttt{LevyItoDecomposition},
ensuring that the downstream abstract framework (truncation convergence, chaos
orthogonality, predictable module structure) receives its inputs as derived
theorems rather than independent axioms. This architectural choice minimizes
the number of primitive assumptions in the formalization.
\end{defn}

\section{Lévy--Itô Decomposition}

\begin{defn}[Lévy--Itô decomposition]\label{def:levy_ito}
\lean{LevyItoDecomposition}
\leanok
\uses{def:toAxioms}
The Lévy--Itô decomposition splits the stable process into small and large jumps:
$L_T = X_T^{\mathrm{small}} + Y_T^{\mathrm{large}}$, where
$X_T^{\mathrm{small}} = \int_0^T \int_{|z|\leq 1} z\,\widetilde{N}(ds,dz)$ captures
compensated small jumps and
$Y_T^{\mathrm{large}} = \sum_{|\Delta L_s|>1} \Delta L_s$ captures large jumps.
In the formalization, this structure packages the compensated integral construction
together with its derived properties (truncation convergence, centering of truncated
integrals, chaos orthogonality, and predictable module structure) into a single
interface that the abstract Banach energy space framework can consume.
\end{defn}

\begin{theorem}[Truncation convergence]\label{thm:trunc_conv}
\lean{LevyItoDecomposition.trunc_convergence}
\leanok
\uses{def:levy_ito}
The truncated pairings converge:
$\mathbb{E}[F \cdot \delta^M(u)] \to \mathbb{E}[F \cdot \delta(u)]$ as $M \to \infty$.
This is derived from the $I_1$-convergence of the constructed compensated integral:
as the upper truncation $M$ increases, the truncated integral
$\delta_L^M(u) = \int\!\int_{|z|\leq M} u_t\,z\,\widetilde{N}(dt,dz)$ converges to
the full integral $\delta_L(u)$ in $L^p$, and Hölder's inequality transfers this to
convergence of the pairing with any $F \in L^q$. This convergence is essential for
the chaos characterization (Theorem~C): it allows passage from the $L^2$ isometry
(valid for truncated integrals) to statements about the full integral $\delta_L$.
\end{theorem}

\begin{theorem}[$\delta^M$ centered]\label{thm:delta_centered}
\lean{LevyItoDecomposition.δ_trunc_centered}
\leanok
\uses{def:levy_ito}
The truncated integral is centered: $\mathbb{E}[\delta^M(u)] = 0$ for all $M > 0$
and all integrands $u$. This is derived from the centering of the constructed
compensated integral $I_1^M$: the truncated integral $\delta^M(u) = I_1^M(u \cdot z_M)$
inherits centering from $I_1^M$. Centering of truncated integrals is used in the
chaos orthogonality argument: it ensures that constant functionals ($F = c$) pair
trivially with $\delta^M(u)$, isolating the non-constant chaos components.
\end{theorem}

\begin{theorem}[Chaos orthogonality]\label{thm:chaos_orth}
\lean{LevyItoDecomposition.chaos_orthogonality_from_kernel_vanishing}
\leanok
\uses{def:levy_ito}
If the truncated chaos pairing $\Phi_M(F) = \int_{|z|\leq M} f_1(s,z)\,z\,\nu_\gamma(dz)$
vanishes (i.e., $\Phi_M \equiv 0$), then $\mathbb{E}[F \cdot \delta^M(u)] = 0$ for all
deterministic integrands $u$. This is derived from the chaos identity
$\mathbb{E}[I_1(f_1) \cdot I_1(u \cdot z_M)] = \int_0^T u(s)\,\Phi_M(s)\,ds$
(the $L^2$ isometry for Poisson integrals) together with \texttt{zero\_mul}.
The vanishing condition $\Phi_M \equiv 0$ holds whenever the first-chaos kernel $f_1$
is even in $z$, since the factor $z$ makes the integrand odd and the symmetric
measure $\nu_\gamma$ annihilates odd integrands. This is the mechanism by which
even functionals (like the centered jump count) are orthogonal to all one-parameter
stochastic integrals.
\end{theorem}

\begin{theorem}[Predictable module]\label{thm:pred_module}
\lean{LevyItoDecomposition.predictable_module_from_linearity}
\leanok
\uses{def:levy_ito}
The truncated pairing respects scalar multiplication:
$\mathbb{E}[(cF) \cdot \delta^M(u)] = c \cdot \mathbb{E}[F \cdot \delta^M(u)]$
for any constant $c$. This is derived from the linearity of the truncated integral
and the linearity of expectation. The predictable module property ensures that the
derivative $D_L$ respects the module structure of $L^q(\Omega)$ over $\mathbb{R}$:
$D_L(cF) = c \cdot D_L(F)$. While seemingly elementary, this property must be
verified for each layer of the construction (finite sums, then limits) to ensure
the abstract framework's linearity axioms are satisfied.
\end{theorem}

\section{Poisson Chaos}

\begin{defn}[Poisson inner product]\label{def:poisson_inner}
\lean{poisson_inner}
\leanok
The Poisson inner product on functions $f, g : \mathbb{N} \to \mathbb{R}$ is defined as
$\langle f, g \rangle_r = \sum_{n=0}^{\infty} f(n)\,g(n)\,P_r(n)$,
where $P_r(n) = e^{-r} r^n / n!$ is the Poisson PMF with parameter $r > 0$.
This is the $L^2$ inner product on the canonical Poisson space
$(\mathbb{N}, \mathrm{poissonMeasure}(r))$. The inner product provides the
concrete realization of the first Poisson chaos: functions of the jump count
are expanded in an orthogonal basis with respect to this pairing, and the
first-chaos projection is governed by the centered variable $\chi(n) = n - r$.
\end{defn}

\begin{defn}[Centered Poisson variable]\label{def:poisson_chi}
\lean{poisson_chi}
\leanok
The centered Poisson variable $\chi : \mathbb{N} \to \mathbb{R}$ is defined by
$\chi(n) = n - r$, where $r$ is the Poisson parameter. This is the canonical
first-chaos element on the Poisson space: it spans the one-dimensional first chaos
$\mathcal{C}_1$ of the large-jump count. Any function $f$ on $\mathbb{N}$ with
$\langle f, \chi \rangle_r = 0$ is orthogonal to the entire first chaos,
meaning it has no ``linear in jump count'' component. The centered witness
$G_0 - \mathbb{E}[G_0]$ for the representability obstruction is precisely
$\chi$ evaluated on the canonical Poisson space with $r = \lambda T$.
\end{defn}

\begin{theorem}[$\chi$ centered]\label{thm:chi_centered}
\lean{poisson_chi_centered}
\leanok
\uses{def:poisson_chi, thm:poisson_mean}
The centered Poisson variable has zero expectation:
$\mathbb{E}[\chi] = \sum_n (n - r) P_r(n) = \mathbb{E}[N] - r = r - r = 0$.
This follows directly from the Poisson mean identity $\mathbb{E}[N] = r$.
Centering of $\chi$ is required for the obstruction argument: the witness
$G = G_0 - \mathbb{E}[G_0]$ must be centered to apply the representability
obstruction theorem, which requires $\mathbb{E}[G] = 0$.
\end{theorem}

\begin{theorem}[$\chi$ norm squared]\label{thm:chi_norm}
\lean{poisson_chi_norm_sq}
\leanok
\uses{def:poisson_chi, thm:poisson_variance}
The $L^2$ norm of the centered Poisson variable equals the parameter:
$\|\chi\|^2 = \sum_n (n-r)^2 P_r(n) = \mathrm{Var}(N) = r$.
This follows from the Poisson variance identity. Since $r = \lambda T > 0$
(as $\lambda > 0$ and $T > 0$), we have $\|\chi\|^2 > 0$, which implies
$\chi \neq 0$ in $L^2$. This is the quantitative content behind the statement
that the centered witness is nonzero: it has positive $L^2$ norm equal to $\lambda T$.
\end{theorem}

\begin{theorem}[First chaos orthogonal]\label{thm:first_chaos}
\lean{first_chaos_orthogonal}
\leanok
\uses{def:poisson_inner, def:poisson_chi}
If $\langle f, \chi \rangle_r = 0$ then $\langle f, a\chi \rangle_r = 0$ for any
scalar $a \in \mathbb{R}$. The proof is immediate from bilinearity:
$\langle f, a\chi \rangle = a \langle f, \chi \rangle = a \cdot 0 = 0$,
proved via \texttt{tsum\_mul\_left}. This result shows that orthogonality to $\chi$
implies orthogonality to the entire one-dimensional first chaos $\mathrm{span}(\chi)$.
It is the concrete realization of the chaos orthogonality principle: a functional
$F$ with $\langle F, \chi \rangle = 0$ has zero first-chaos coefficient and is therefore
invisible to the one-parameter derivative $D_L$.
\end{theorem}

\begin{theorem}[Chaos orthogonality concrete]\label{thm:chaos_concrete}
\lean{chaos_orthogonality_concrete}
\leanok
\uses{thm:first_chaos}
Zero first-chaos coefficient implies orthogonality to the entire first chaos:
if $f$ has $\langle f, \chi \rangle_r = 0$, then $f$ is orthogonal to every element
of $\mathrm{span}(\chi)$. This is the concrete instantiation of the abstract chaos
orthogonality theorem on the canonical Poisson space. It provides the bridge between
the Poisson-space computation (working with functions on $\mathbb{N}$) and the abstract
Banach framework (working with $L^q$ functionals on $\Omega$): a functional whose
first-chaos coefficient vanishes lies in $\ker(D_L)$, confirming that it cannot be
represented as a one-parameter stochastic integral $\delta_L(u)$.
\end{theorem}

\section{$L^2$ Estimates}

\begin{theorem}[$L^2$ Cauchy estimate]\label{thm:l2_cauchy}
\lean{l2_cauchy_vanishes}
\leanok
The $L^2$ Cauchy estimate for the compensated integral approximation is
$\|I_1^{\varepsilon_1,M}(h) - I_1^{\varepsilon_2,M}(h)\|_2^2
\leq 2c_\gamma \varepsilon^{2-\gamma}/(2-\gamma) \to 0$ as $\varepsilon \to 0^+$.
This bound comes from the $L^2$ isometry applied to the difference of truncated
integrals, which involves $\int_{\varepsilon_1}^{\varepsilon_2} z^2\,\nu_\gamma(dz)
= 2c_\gamma(\varepsilon_2^{2-\gamma} - \varepsilon_1^{2-\gamma})/(2-\gamma)$.
Since $2 - \gamma > 0$, the right-hand side vanishes as both $\varepsilon_i \to 0$.
This estimate is the technical heart of the $\varepsilon \to 0$ construction:
it proves that the net of approximations is Cauchy in $L^2$, guaranteeing
convergence to the compensated integral $I_1^M(h)$.
\end{theorem}

\begin{theorem}[rpow tends to zero]\label{thm:rpow_zero}
\lean{rpow_tendsto_zero_at_zero_right}
\leanok
The power function $\varepsilon^\alpha \to 0$ as $\varepsilon \to 0^+$ for any $\alpha > 0$.
This is a basic real-analysis fact used in the $L^2$ Cauchy estimate: the bound
$2c_\gamma \varepsilon^{2-\gamma}/(2-\gamma)$ involves the exponent $2 - \gamma > 0$,
so $\varepsilon^{2-\gamma} \to 0$ as $\varepsilon \to 0^+$. The result is also used
in the moment divergence/convergence analysis and whenever truncation parameters
tend to boundary values. In the Lean formalization, this is proved using Mathlib's
\texttt{Filter.Tendsto} API for limits at the boundary of the positive reals.
\end{theorem}

\begin{theorem}[Truncated moment bounded]\label{thm:trunc_bounded}
\lean{truncated_moment_bounded}
\leanok
The truncated second moment $\int_{|z|\leq M} z^2\,\nu_\gamma(dz) = 2c_\gamma M^{2-\gamma}/(2-\gamma)$
is positive and finite for any fixed $M > 0$. This explicit formula follows from
the power-law computation $\int_0^M z^{2-1-\gamma}\,dz = M^{2-\gamma}/(2-\gamma)$
with the factor of $2$ from symmetry. The bound provides the isometry constant for
the truncated compensated integral $\delta_L^M$: the $L^2$ isometry reads
$\mathbb{E}[|\delta_L^M(u)|^2] = \|u\|_2^2 \cdot \int_{|z|\leq M} z^2\,\nu_\gamma(dz)$.
Finiteness ensures the isometry is well-defined; positivity ensures it is nondegenerate.
\end{theorem}

\section{Compound Poisson Path}

\begin{defn}[Compound Poisson path]\label{def:cp_path}
\lean{compound_poisson_path}
\leanok
\uses{def:cpi}
The compound Poisson path is the concrete stochastic process
$Y_t = \sum_{s_i \leq t} z_i$, where $\{(s_i, z_i)\}$ are the jump times and sizes
of a compound Poisson process. This is a pure-jump càdlàg process that increases by
$z_i$ at each jump time $s_i$. The path is defined as a deterministic function of the
jump data (times and sizes), making it fully constructive. It serves as the concrete
realization of both the small-jump component (with $|z_i| \leq M$, after compensation)
and the large-jump component (with $|z_i| > 1$) in the Lévy--Itô decomposition.
\end{defn}

\begin{defn}[Stable large jump path]\label{def:sljp}
\lean{StableLargeJumpPath}
\leanok
\uses{def:cpi}
The stable large jump path is a compound Poisson path restricted to jumps with
$|z_i| > 1$: $Y_t^{\mathrm{large}} = \sum_{s_i \leq t, |z_i| > 1} z_i$.
Since the Lévy measure $\nu_\gamma$ has finite mass on $\{|z| > 1\}$
($\nu_\gamma(\{|z|>1\}) = 2c_\gamma/\gamma < \infty$), there are a.s.\
finitely many large jumps in $[0,T]$. The number of large jumps
$N_T = \#\{s \leq T : |\Delta L_s| > 1\}$ is a Poisson random variable
with parameter $\lambda T = 2c_\gamma T/\gamma$, and the centered version
$N_T - \mathbb{E}[N_T]$ serves as the obstruction witness.
\end{defn}

\begin{theorem}[Path starts at zero]\label{thm:path_zero}
\lean{StableLargeJumpPath.path_zero}
\leanok
\uses{def:cp_path}
The compound Poisson path satisfies $Y_0 = 0$: at time $t = 0$, the sum
$\sum_{s_i \leq 0} z_i$ is empty, giving zero. This is a standard initial condition
for Lévy processes and is verified by showing that the sum over an empty index set
vanishes. The result ensures that the Itô formula $f(Y_T) - f(Y_0) = f(Y_T) - f(0)$
starts from the correct initial value, which is used in the telescoping proof
of the compound Poisson Itô formula.
\end{theorem}

\begin{theorem}[Path terminal value]\label{thm:path_terminal}
\lean{StableLargeJumpPath.path_terminal}
\leanok
\uses{def:cp_path}
The terminal value of the compound Poisson path equals the sum of all jumps:
$Y_T = \sum_{i} z_i$. This is immediate from the definition $Y_t = \sum_{s_i \leq t} z_i$
evaluated at $t = T$, since all jump times satisfy $s_i \leq T$ by construction.
The terminal value identity connects the path-level description (a function of time)
to the random-variable description (a sum of i.i.d.\ jump sizes), which is needed
when applying the Itô formula to terminal functionals $F = f(L_T)$.
\end{theorem}


\chapter{Banach Energy Space Framework}

\section{Abstract Framework}

\begin{defn}[Banach energy space]\label{def:banach_space}
\lean{BanachEnergySpace}
\leanok
The abstract Banach energy space packages the data needed for the operator factorization
framework: a Banach integrand space $H_X$, its dual $H_X^*$, function spaces
$L^p(\Omega)$ and $L^q(\Omega)$ with conjugate exponents, a bounded stochastic integral
(divergence) $\delta : H_X \to L^p(\Omega)$, bilinear duality pairings, an expectation
functional, and a \texttt{mk\_dual} constructor that packages bounded linear functionals
into elements of $H_X^*$. Unlike the Hilbert framework of the companion paper, this
structure does \emph{not} assume a Riesz identification $H_X^* \cong H_X$: the dual space
is genuinely distinct from the integrand space when $p \neq 2$.
The framework isolates the stochastic-analysis content (integration, expectation)
from the Banach-space algebra (duality, linearity), enabling a clean separation of
concerns in the formalization.
\end{defn}

\begin{defn}[Centering]\label{def:center}
\lean{BanachEnergySpace.center}
\leanok
\uses{def:banach_space}
The centering operator maps a random variable $G_0$ to its centered version:
$\mathrm{center}(G_0) = G_0 - \mathrm{const}(\mathbb{E}[G_0])$, where
$\mathrm{const}(c)$ denotes the constant random variable taking value $c$.
Centering removes the mean while preserving all other distributional properties
(variance, symmetry, chaos decomposition beyond the zeroth chaos).
In the obstruction argument, centering is applied to the large-jump count $G_0$
to produce a centered, even, nonzero witness $G = G_0 - \mathbb{E}[G_0]$ that
lies in $\ker(D_L) \setminus \{0\}$.
\end{defn}

\begin{theorem}[Center is centered]\label{thm:center_centered}
\lean{BanachEnergySpace.center_centered}
\leanok
\uses{def:center}
The centered version of any random variable has zero expectation:
$\mathbb{E}[\mathrm{center}(G_0)] = \mathbb{E}[G_0 - \mathrm{const}(\mathbb{E}[G_0])]
= \mathbb{E}[G_0] - \mathbb{E}[G_0] = 0$.
This is immediate from linearity of expectation and the property
$\mathbb{E}[\mathrm{const}(c)] = c$. The centered expectation is one of the three
hypotheses of the representability obstruction theorem (the others being evenness
and nonzeroness): only centered functionals can potentially lie in $\mathrm{im}(\delta_L)$
(since all stochastic integrals are centered by (H1)), so centering is a necessary
preparation for testing representability.
\end{theorem}

\section{The Derivative $D$}

\begin{defn}[Operator-covariant derivative]\label{def:D}
\lean{BanachEnergySpace.D}
\leanok
\uses{def:banach_space}
The operator-covariant derivative is \emph{constructed} (not axiomatized) as the
Banach dual of the stochastic integral:
$D(F) := \mathrm{mk\_dual}(u \mapsto \mathbb{E}[F \cdot \delta(u)])$.
For each $F \in L^q(\Omega)$, the map $u \mapsto \mathbb{E}[F \cdot \delta(u)]$ defines
a bounded linear functional on $H_X$ (by Hölder's inequality:
$|\mathbb{E}[F \cdot \delta(u)]| \leq \|F\|_q \cdot \|\delta(u)\|_p
\leq \|F\|_q \cdot C_\delta \|u\|_{H_X}$), so $D(F) \in H_X^*$.
This construction is the Banach-space generalization of the Hilbert derivative
$D_X = \delta_X^*$, with the crucial difference that $D(F)$ lives in $H_X^*$
(the dual space), not in $H_X$ itself. The defining identity
$\langle DF, u\rangle = \mathbb{E}[F \cdot \delta(u)]$ is then a \emph{theorem},
not a definition.
\end{defn}

\begin{theorem}[$D$ defining identity]\label{thm:D_def}
\lean{BanachEnergySpace.D_def}
\leanok
\uses{def:D}
The defining identity of the operator-covariant derivative is
$\langle DF, u \rangle = \mathbb{E}[F \cdot \delta(u)]$ for all $F \in L^q(\Omega)$
and $u \in H_X$. This identity is proved (in one line) from the construction of $D$
via \texttt{mk\_dual}: the pairing of the packaged functional with $u$ unfolds to the
expectation. This identity is the starting point for all downstream results: linearity
of $D$, vanishing on constants, the kernel characterization, and the Leibniz defect
formula all follow from manipulating this single pairing identity.
\end{theorem}

\begin{theorem}[$D$ linear]\label{thm:D_linear}
\lean{BanachEnergySpace.D_linear}
\leanok
\uses{thm:D_def}
The derivative $D$ is linear: $D(\alpha F + G) = \alpha \cdot DF + DG$ for all
$\alpha \in \mathbb{R}$ and $F, G \in L^q(\Omega)$. The proof evaluates both sides
on an arbitrary test integrand $u$ using the defining identity:
$\langle D(\alpha F + G), u \rangle = \mathbb{E}[(\alpha F + G) \cdot \delta(u)]
= \alpha \mathbb{E}[F \cdot \delta(u)] + \mathbb{E}[G \cdot \delta(u)]
= \alpha \langle DF, u \rangle + \langle DG, u \rangle$.
Nondegeneracy of the pairing then gives equality in $H_X^*$.
Linearity of $D$ is essential for interpreting $\ker(D)$ as a linear subspace and
$D$ as a bounded linear operator $L^q(\Omega) \to H_X^*$.
\end{theorem}

\begin{theorem}[$D$ on constants]\label{thm:D_const}
\lean{BanachEnergySpace.D_const}
\leanok
\uses{thm:D_def}
The derivative of a constant random variable vanishes: $D(\mathrm{const}(c)) = 0$
for all $c \in \mathbb{R}$. The proof uses the defining identity and (H1):
$\langle D(\mathrm{const}(c)), u \rangle = \mathbb{E}[c \cdot \delta(u)]
= c \cdot \mathbb{E}[\delta(u)] = c \cdot 0 = 0$ for all $u$, so
$D(\mathrm{const}(c)) = 0$ by nondegeneracy. This result, combined with linearity,
shows that $D$ factors through the centering map: $D(F) = D(F - \mathbb{E}[F])$,
so the derivative depends only on the fluctuation $F - \mathbb{E}[F]$ and is
blind to the mean.
\end{theorem}

\begin{theorem}[$\ker(D)$ characterization]\label{thm:ker_D}
\lean{BanachEnergySpace.ker_D_iff_annihilates_im_delta}
\leanok
\uses{thm:D_def}
The kernel of $D$ is characterized as the annihilator of $\mathrm{im}(\delta)$:
$DF = 0$ if and only if $\mathbb{E}[F \cdot \delta(u)] = 0$ for all $u \in H_X$.
The forward direction is immediate from the defining identity; the reverse uses
nondegeneracy of the $H_X^*$--$H_X$ pairing. This characterization reveals the
operator-theoretic content of the kernel: $F \in \ker(D)$ means $F$ is
$L^q$--$L^p$-orthogonal to every stochastic integral, i.e., $F$ carries no
information detectable by the one-parameter integral. For stable processes,
$\ker(D_L)$ is large (containing all even functionals and all higher Poisson
chaoses), which is the source of the representability obstruction.
\end{theorem}

\section{Fluctuation Factorization (Theorem A)}

\begin{theorem}[Fluctuation factorization]\label{thm:factorization}
\lean{BanachEnergySpace.fluctuation_factorization}
\leanok
\uses{def:D, def:banach_space}
\textbf{(Theorem~A.)} Under the representation hypothesis (H3), every centered
functional factors through the stochastic integral:
$(F - \mathbb{E}[F]) = \delta(\delta^{-1}(F - \mathbb{E}[F]))$ for all
$F \in L^p(\Omega) \cap L^q(\Omega)$.
The proof is direct: (H3) gives $F - \mathbb{E}[F] = \delta(v)$ for some predictable
$v$, and (H2) ensures $v$ is unique, so $v = \delta^{-1}(F - \mathbb{E}[F])$.
For symmetric $\gamma$-stable processes, (H3) fails globally (the one-parameter integral
cannot represent even functionals), but the factorization still holds on the closed
subspace $\mathrm{im}(\delta_L) \subsetneq L^p_0(\Omega)$, characterizing exactly
which functionals admit one-parameter representation.
\end{theorem}

\begin{theorem}[Duality recovery]\label{thm:duality_recovery}
\lean{BanachEnergySpace.duality_recovery}
\leanok
\uses{thm:factorization, thm:D_def}
The duality recovery identity states that
$\langle DF, u \rangle = \mathbb{E}[\delta(v) \cdot \delta(u)]$
where $v = \delta^{-1}(F - \mathbb{E}[F])$ is the representing integrand.
This shows that the derivative $D$ \emph{recovers} the integrand from the
functional: the pairing $\langle DF, u \rangle$ encodes the correlation structure
of $F$ with all stochastic integrals, and under (H3) this uniquely determines
the representing integrand $v$. The proof uses (H1) to eliminate the constant
$\mathbb{E}[F]$: $\langle DF, u \rangle = \mathbb{E}[F \cdot \delta(u)]
= \mathbb{E}[(F - \mathbb{E}[F]) \cdot \delta(u)] = \mathbb{E}[\delta(v) \cdot \delta(u)]$.
\end{theorem}

\section{Representability Obstruction}

\begin{defn}[Distributional symmetry]\label{def:dist_sym}
\lean{BanachEnergySpace.DistributionalSymmetry}
\leanok
\uses{def:banach_space}
Distributional symmetry is the property $L \stackrel{d}{=} -L$: the process $L$ and
its negation $-L$ have the same law. For symmetric $\gamma$-stable processes, this
holds because $\nu_\gamma$ is symmetric ($\nu_\gamma(-dz) = \nu_\gamma(dz)$) and the
process has no drift. Distributional symmetry induces a negation map on $L^q(\Omega)$
that preserves functionals depending only on $|\Delta L_s|$ (even functionals)
while flipping the sign of stochastic integrals $\delta_L(u) = \int u_t\,dL_t$
(odd functionals). This parity mismatch is the mechanism behind the even annihilation
lemma and, ultimately, the representability obstruction.
\end{defn}

\begin{theorem}[Even annihilates $\mathrm{im}(\delta)$]\label{thm:even_annihilates}
\lean{BanachEnergySpace.even_annihilates_im_delta}
\leanok
\uses{def:dist_sym, thm:D_def}
If $G$ is even (i.e., $G \circ \mathrm{neg} = G$, invariant under the distributional
symmetry), then $\mathbb{E}[G \cdot \delta(u)] = 0$ for all $u \in H_X$.
The proof uses the distributional symmetry $L \stackrel{d}{=} -L$: under the sign flip,
$G$ is invariant while $\delta(u)$ changes sign, so
$\mathbb{E}[G \cdot \delta(u)] = \mathbb{E}[G \cdot (-\delta(u))]
= -\mathbb{E}[G \cdot \delta(u)]$, forcing the expectation to zero.
This lemma is the heart of the obstruction: it shows that all even functionals
lie in $\ker(D_L)$, creating a large kernel that prevents (H3) from holding.
\end{theorem}

\begin{theorem}[Representability obstruction]\label{thm:obstruction}
\lean{BanachEnergySpace.representability_obstruction}
\leanok
\uses{thm:even_annihilates, thm:center_centered}
\textbf{(Proposition 5.1.)} If $G \in L^q(\Omega)$ is centered ($\mathbb{E}[G] = 0$),
even ($G \circ \mathrm{neg} = G$), and nonzero, then $G \notin \mathrm{im}(\delta_L)$.
The proof is by contradiction: if $G = \delta(u)$ for some $u$, then the self-pairing
$\mathbb{E}[G \cdot G] = \mathbb{E}[G \cdot \delta(u)] = 0$ by the even annihilation
lemma, which forces $\|G\|_2 = 0$, hence $G = 0$---contradicting nonzeroness.
This establishes that $\mathrm{im}(\delta_L) \subsetneq L^p_0(\Omega)$: hypothesis
(H3) fails for stable processes. The obstruction is structural, not technical: it
reflects the mismatch between one-parameter integrands (linear in jump size $z$)
and the full Poisson chaos (which includes nonlinear dependence on $z$).
\end{theorem}

\section{Leibniz Defect (Theorem B)}

\begin{defn}[Leibniz defect]\label{def:leibniz}
\lean{BanachEnergySpace.leibniz_defect}
\leanok
\uses{def:D}
The Leibniz defect measures the failure of the product rule for the derivative $D$:
$\Gamma(F,G) := D(FG) - F \cdot DG - G \cdot DF$, viewed as an element of $H_X^*$.
When $D$ is a derivation (satisfies the Leibniz rule), $\Gamma = 0$; nonzero $\Gamma$
is the algebraic signature of stochastic calculus. In the Brownian (Hilbert) setting,
$\Gamma$ encodes the carré du champ $\int f'(B_t) g'(B_t) u_t\,dt$. In the stable
(Banach) setting, $\Gamma$ captures the full nonlinear jump residual, involving all
orders of $f$ and $g$ via the product jump defect $\Delta_s(f,g)$. The Leibniz defect
is a standalone duality object: its computation requires only $\langle D(FG), u \rangle
= \mathbb{E}[FG \cdot \delta(u)]$, the Itô formula, and Hölder's inequality---not
the representation hypothesis (H3).
\end{defn}

\begin{theorem}[Product rule with jump defect]\label{thm:product_rule}
\lean{BanachEnergySpace.product_rule_with_jump_defect}
\leanok
\uses{def:leibniz}
\textbf{(Theorem~B.)} For $f, g \in C_b^2(\mathbb{R})$, $F = f(L_T)$, $G = g(L_T)$,
the Leibniz defect is given by the product jump defects:
$\langle \Gamma(F,G), u \rangle = \mathbb{E}[\sum_{0 < s \leq T} \Delta_s(f,g) \cdot \delta(u)]$,
where $\Delta_s(f,g) = (fg)(L_s) - (fg)(L_{s-}) - [f'g + fg'](L_{s-}) \Delta L_s$.
The proof applies the Lévy--Itô formula to $h = fg$, multiplies by $\delta_L(u)$,
takes expectations, and identifies the Leibniz part with the predictable integrand
$(f'g + fg')(L_{t-})$. The jump defect sum converges a.s.\ because small jumps
contribute $O(|\Delta L_s|^2)$ (summable since $\gamma < 2$) and large jumps are
finitely many. This theorem demonstrates that the one-parameter derivative $D_L$
\emph{detects} the full nonlinear jump structure through its algebraic failure to
satisfy the product rule, even though it cannot \emph{represent} all centered functionals.
\end{theorem}

\begin{theorem}[Defect vanishes without jumps]\label{thm:defect_vanishes}
\lean{BanachEnergySpace.defect_vanishes_without_jumps}
\leanok
\uses{def:leibniz}
In the absence of jumps (i.e., when the product jump defect
$\sum_s \Delta_s(f,g) = 0$), the Leibniz defect vanishes: $\Gamma(F,G) = 0$.
This recovers the classical product rule $D(FG) = F \cdot DG + G \cdot DF$
that holds for continuous processes (e.g., Brownian motion). The result confirms
that the Leibniz defect is a \emph{jump phenomenon}: it arises exclusively from
the discontinuities of the driving process. In the Hilbert bridge
(\texttt{HilbertEnergyData}), this vanishing condition is the hypothesis that
specializes the Banach framework back to the Hilbert theory.
\end{theorem}

\section{Chaos Characterization (Theorem C)}

\begin{theorem}[Chaos kernel characterization]\label{thm:chaos_char}
\lean{BanachEnergySpace.chaos_kernel_characterization}
\leanok
\uses{thm:D_def}
\textbf{(Theorem~C.)} The derivative $D_L F$ vanishes on deterministic integrands
if and only if the truncated pairings
$\Phi_M(s) = \int_{|z|\leq M} f_1(s,z)\,z\,\nu_\gamma(dz) \to 0$ in
$L^q([0,T])$ as $M \to \infty$, where $f_1$ is the first-chaos kernel of $F$.
The proof uses the $L^2$ isometry for truncated Poisson integrals to reduce
$\mathbb{E}[F \cdot \delta_L^M(u)] = \int_0^T u(s)\,\Phi_M(s)\,ds$ (chaos
orthogonality eliminates all but the first chaos), then passes to $M \to \infty$
via truncation convergence. The vanishing condition $\Phi_M \to 0$ holds automatically
for even kernels ($f_1(s,-z) = f_1(s,z)$, since the integrand $f_1 \cdot z$ is odd
and $\nu_\gamma$ is symmetric) and for all higher chaoses ($n \geq 2$, by orthogonality).
The annihilator of $\mathrm{im}(\delta_L)$ is therefore infinite-dimensional,
containing all higher chaoses and an infinite-dimensional subspace of the first chaos.
\end{theorem}

\section{Stable Lévy Instantiation}

\begin{defn}[Stable noise data]\label{def:stable_noise}
\lean{StableNoiseData}
\leanok
\uses{def:banach_space, def:dist_sym}
The stable noise data packages the concrete parameters of a symmetric $\gamma$-stable
Lévy process into the abstract Banach energy space framework: the stability index
$\gamma \in (1,2)$, the distributional symmetry $L \stackrel{d}{=} -L$, a jump count
random variable $G_0 = \#\{|\Delta L_s| > 1, s \leq T\}$, the large-jump rate
$\lambda = 2c_\gamma/\gamma > 0$, and the terminal time $T > 0$.
The key structural properties are that $G_0$ factors through absolute values
of jumps (making $\mathrm{center}(G_0)$ even) and that $\lambda T > 0$
(making $\mathrm{center}(G_0)$ nonzero via positive Poisson variance).
These two properties, combined with centering, feed the three hypotheses of the
abstract representability obstruction theorem.
\end{defn}

\begin{theorem}[Variance positive]\label{thm:var_pos}
\lean{jump_count_variance_pos}
\leanok
\uses{def:stable_noise}
The variance of the large-jump count is strictly positive:
$\mathrm{Var}(G_0) = \lambda T > 0$, since $\lambda = 2c_\gamma/\gamma > 0$
(from positivity of $c_\gamma$ and $\gamma$) and $T > 0$ (positive terminal time).
The proof in Lean uses \texttt{mul\_pos} to combine the two positivity hypotheses.
Positive variance is the quantitative statement that the centered witness
$G_0 - \mathbb{E}[G_0]$ is nonzero: $\|G_0 - \mathbb{E}[G_0]\|_2^2 = \lambda T > 0$
implies $G_0 - \mathbb{E}[G_0] \neq 0$ in $L^2$.
\end{theorem}

\begin{theorem}[Centered witness even]\label{thm:witness_even}
\lean{jump_count_center_even}
\leanok
\uses{def:stable_noise, thm:center_centered}
The centered large-jump count $\mathrm{center}(G_0) = G_0 - \mathbb{E}[G_0]$ is even
under the distributional symmetry. The proof has two parts: (1) $G_0$ is even because
it counts jumps $|\Delta L_s| > 1$, which depends only on absolute jump sizes and is
invariant under $L \mapsto -L$ (formalized via the \texttt{jump\_count\_factors\_through\_abs}
axiom and \texttt{const\_even}); (2) centering preserves evenness because
$\mathbb{E}[G_0]$ is a constant (hence even), and the difference of even functions
is even (via \texttt{neg\_sub}). Evenness of the centered witness is one of the three
hypotheses of the representability obstruction theorem.
\end{theorem}

\begin{theorem}[Centered witness nonzero]\label{thm:witness_nonzero}
\lean{jump_count_center_nonzero}
\leanok
\uses{thm:var_pos}
The centered large-jump count is nonzero: $\mathrm{center}(G_0) \neq 0$ in $L^2$.
This follows from the positive variance: $\|\mathrm{center}(G_0)\|_2^2
= \mathrm{Var}(G_0) = \lambda T > 0$, so the $L^2$ norm is strictly positive,
which implies the function is nonzero. The derivation chain in Lean is:
$\lambda > 0$ and $T > 0$ $\Rightarrow$ $\lambda T > 0$ (via \texttt{mul\_pos})
$\Rightarrow$ $\mathrm{Var}(G_0) > 0$ $\Rightarrow$ $\mathrm{center}(G_0) \neq 0$.
This is the final ingredient needed for the stable obstruction: combined with
centering and evenness, it triggers the abstract representability obstruction theorem.
\end{theorem}

\begin{theorem}[Stable obstruction]\label{thm:stable_obstruction}
\lean{stable_levy_obstruction}
\leanok
\uses{thm:obstruction, thm:witness_even, thm:witness_nonzero, thm:center_centered}
The centered large-jump count is not representable as a one-parameter stochastic integral:
$G_0 - \mathbb{E}[G_0] \notin \mathrm{im}(\delta_L)$.
This is the concrete instantiation of the abstract representability obstruction
applied to the stable Lévy setting. The proof assembles three derived properties:
(1)~$\mathrm{center}(G_0)$ is centered (from \texttt{center\_centered}),
(2)~$\mathrm{center}(G_0)$ is even (from \texttt{jump\_count\_center\_even}),
(3)~$\mathrm{center}(G_0) \neq 0$ (from \texttt{jump\_count\_center\_nonzero}).
The abstract obstruction theorem then concludes $G \notin \mathrm{im}(\delta_L)$.
This is the paper's central negative result: it shows that hypothesis (H3) fails
for symmetric $\gamma$-stable processes, proving that the one-parameter factorization
$(\mathrm{Id} - \mathbb{E}) = \delta_L \circ \delta_L^{-1} \circ (\mathrm{Id} - \mathbb{E})$
does not hold on all of $L^p_0(\Omega)$.
\end{theorem}

\section{Hilbert Bridge}

\begin{theorem}[Leibniz defect vanishes]\label{thm:hilbert_bridge}
\lean{HilbertEnergyData.leibniz_defect_vanishes}
\leanok
\uses{def:leibniz}
In the Hilbert case (continuous processes with no jumps), the Leibniz defect vanishes:
$\Gamma(F,G) = 0$ for all $F, G$. The \texttt{HilbertEnergyData} bridge shows that
every \texttt{BanachEnergySpace} axiom is satisfied when the product jump defect
$\sum_s \Delta_s(f,g)$ is identically zero, recovering the Hilbert theory as a special
case. This confirms that the Banach framework is a strict generalization: the Hilbert
operator factorization of the companion paper is the specialization where $p = q = 2$,
$H_X^* \cong H_X$ via Riesz, and $\Gamma = 0$ because there are no jumps.
The bridge validates the design of the Banach framework by ensuring backward
compatibility with the previously verified Hilbert formalization.
\end{theorem}
